\documentclass{article}
\usepackage{tcee-repaso}
\tcetema[La inflación]{3.A.41}{La inflación: causas y efectos sobre la eficiencia económica y el bienestar. Hiperinflación y deflación.}
\tcefecha{05/11/2026}
\begin{document}
\cabecerarepaso
\begin{columnas}

% ==========================================================================
\bloque{Esquema del cante}
\esquemacante{%
\apartadocante{Introducción: «quien juega con fuego, se quema»}{2,5}
\apartadocante{I. Causas de la inflación}{13}
\subcante{1. Debate tradicional: clásicos y keynesianos}{1,5}
\subcante{2. Modelos dinámicos: el triángulo de \autor{Gordon}}{11,5}
\apartadocante{II. Efectos sobre la eficiencia y el bienestar}{8}
\subcante{1. Costes de la anticipada y la no anticipada}{4}
\subcante{2. ¿Existe una tasa de inflación óptima?}{3}
\subcante{3. Inflación y estabilización}{1}
\apartadocante{III. Contextos extremos}{5,5}
\subcante{1. Hiperinflación}{3,5}
\subcante{2. Deflación}{2}
\apartadocante{Conclusión}{1}}
Pegatinas: «echar leña al fuego» (demanda), «la pescadilla que se muerde la cola» (2.ª ronda), «genio y figura, hasta la sepultura» (inercia), «las cuentas claras y el chocolate espeso» (unidad de cuenta), «curarse en salud» (el 2\,\%).

% ==========================================================================
\bloque[2,5]{Introducción}
\begin{itemize}
\item \textbf{Enganche}: «Quien juega con fuego, se quema». Alemania, oct.\ 1923: precios ×$\approx$300 en un mes, se duplicaban cada 3,7 días; déficit pagado con la imprenta (\autor{Hanke}-\autor{Krus}, 2013; \autor{Sargent}, 1982). Un siglo después, la zona del euro temía lo contrario: precios ↓ en 2009, dic.\ 2014-mar.\ 2015, 2016 y ago.-dic.\ 2020. Y vuelve el fuego: oct.\ 2026, zona del euro 4,0\,\% (energía +20,8\,\%), España 4,9\,\%; expectativas a LP ancladas en el 2,0\,\%. ¿\emph{De dónde sale el fuego y cuánto calor conviene}?
\item \textbf{Relevancia}: \emph{teoría}: campo de pruebas de la macro (neutralidad, \autor{Phillips}, expectativas racionales, inconsistencia temporal, rigideces NEK, monetaria-fiscal). \emph{Política}: estabilidad de precios, objetivo principal del SEBC (art.~127.1 TFUE); BCE: 2\,\% simétrico a medio plazo (jul.\ 2021, confirmado jun.\ 2025) \vertema{3.A.36}. \emph{Bienestar}: impuesto que no aprueba ningún parlamento; pesa más en rentas bajas (energía y alimentos).
\item \textbf{Contextualización}: \concepto{inflación}: aumento \emph{sostenido} y \emph{generalizado} del nivel de precios, $\pi_t\equiv(P_t-P_{t-1})/P_{t-1}$; \concepto{deflación}: caída sostenida y generalizada (no basta un mes negativo); \concepto{desinflación}: $\pi$↓ sin dejar de ser $>0$; \concepto{hiperinflación} (\autor{Cagan}, 1956): $>50$\,\% \emph{al mes}. \emph{Medición}: IPC (IAPC en la UE, comparable); deflactor del PIB (lo producido, sin importado); \concepto{subyacente}: sin energía ni alimentos \vertema{3.A.9}; el IPC sobrestima (calidad, productos nuevos, sustitución). \emph{4 etapas}: cuantitativa \vertema{3.A.3}; keynesiana y síntesis (demanda, costes, \autor{Phillips} como menú) \vertema{3.A.5}; expectativas (1968-80s: \autor{Friedman}-\autor{Phelps}, \autor{Lucas}, \autor{Kydland}-\autor{Prescott}, \autor{Sargent}-\autor{Wallace}); nueva síntesis (rigideces + expectativas racionales + bancos centrales independientes; 1999-2020, media zona del euro 1,6\,\%). \emph{Consenso}: a LP $\pi$ sigue al dinero; alta y volátil es mala. \emph{Debate}: qué la mueve a CP y por qué persiste; qué objetivo (\autor{Friedman}: deflación suave; rigideces: 0; límite inferior y salarios: algo $>0$).
\item \textbf{Problemática}: ¿dinero, demanda, costes o inercia? ¿Siempre monetaria? ¿Qué cuesta, depende de que se anticipe, hay tasa óptima? ¿Por qué dañan tanto los extremos y cómo se sale?
\item \textbf{Estructura} (orden del título): \hyperref[rep:I]{I}, causas (de clásicos y keynesianos al triángulo, aplicado a 2021-26); \hyperref[rep:II]{II}, efectos (anticipada y no anticipada, tasa óptima, estabilización); \hyperref[rep:III]{III}, los 2 incendios, demasiado fuego y ninguno (se entienden con las causas y los costes).
\end{itemize}

% ==========================================================================
\bloque[13]{I. Causas de la inflación\label{rep:I}}
\begin{ideaclave}
A largo plazo la inflación es un fenómeno monetario; a corto la explican 3 fuerzas, la demanda, los costes y la inercia, y la política monetaria y la fiscal deciden si el dinero la alimenta y, por tanto, si se queda.
\end{ideaclave}

\sub{I.1. Debate tradicional: clásicos y keynesianos\label{rep:I.1}}
\begin{itemize}
\item \emph{Teoría cuantitativa} (con $Y$ natural y $V$ estable), en tasas \vertema{3.A.3}:
\begin{formula}
\centering $\displaystyle \gamma_M+\gamma_V=\pi+\gamma_Y$
\end{formula}
\item[] \autor{Friedman} (1963): «la inflación es siempre y en todas partes un fenómeno monetario» (solo la produce $M$ creciendo más que $Y$) \vertema{3.A.5}. \emph{Evid.}: 110 países, 1960-90: correlación dinero-inflación 0,925 (base), 0,958 (M1), 0,950 (M2); con $\gamma_Y$, casi nula (\autor{McCandless}-\autor{Weber}, 1995) \vertema{3.A.35}.
\item \emph{Clásicos}: precios flexibles, dinero neutral; el estímulo solo sube $P$ y la inflación desaparece al retirarlo. \emph{Keynesianos}: $W$ rígido a la baja; la demanda mueve antes $Y$ y $L$; inflación de demanda o de costes (puja salarios-beneficios) \vertema{3.A.4}. \emph{Monetaristas}: la sostenida es siempre monetaria; el CP, por expectativas.
\item En los 3, un impulso único sube el \emph{nivel} una vez: para inflación, el impulso se repite o algo lo perpetúa $\Rightarrow$ modelos dinámicos.
\end{itemize}

\sub{I.2. Modelos dinámicos: el triángulo de \autor{Gordon} (1982)\label{rep:I.2}}
\begin{itemize}
\item \emph{La curva de \autor{Phillips}} («probablemente la relación macro más importante», \autor{Akerlof}, 2002): \autor{Phillips} (1958, R.\,U.): $\gamma_W$ vs.\ paro (antes, \autor{Fisher}, 1926); \autor{Samuelson}-\autor{Solow} (1960): a precios, como \emph{menú}; \autor{Friedman} (1968)-\autor{Phelps} (1967): cuenta la $\pi$ \emph{esperada}. \concepto{Aumentada con expectativas} \vertema{3.A.28}:
\begin{formula}
\centering $\displaystyle \pi_t=\pi^e_t-a\,(U_t-U^n)+\varepsilon_t,\quad a>0$
\end{formula}
\item[] $U^n$ tasa natural (NAIRU), $\varepsilon_t$ perturbación de oferta; a LP ($\pi=\pi^e$) vertical en $U^n$.
\item \concepto{Triángulo} (\autor{Gordon}, 1977, 1982; bautizado en 1997): 3 orígenes = 3 términos; con expectativas que miran a los $k$ periodos anteriores:
\funcionclave{\pi_t=\underbrace{\textstyle\sum_{i=1}^{k}a_i\pi_{t-i}}_{\text{inercia}}\underbrace{-\,a\,(U_t-U^n)}_{\text{demanda}}+\underbrace{\varepsilon_t}_{\text{costes}},\ \ \textstyle\sum a_i=1}
\item[] \concepto{De demanda} (\emph{demand-pull}): $U<U^n$ empuja precios; \concepto{de costes} (\emph{cost-push}): perturbación de oferta (energía, materias primas, márgenes, tipo de cambio) con cualquier paro; \concepto{inercial} (\emph{built-in}): sigue por expectativas, indexación, contratos. Pesos suman 1 $\Rightarrow$ aceleracionista: con $U=U^n$ y sin perturbaciones $\pi$ se queda; solo baja si $U>U^n$ (desinflar cuesta). Con histéresis, cuenta $\Delta U$.
\end{itemize}
\pareja{esquema-origen-inflacion}{Triángulo de \autor{Gordon}: 3 vértices = 3 términos}{phillips-aumentada}{\autor{Phillips} aumentada: $A\to B$ (CP), $B\to C$ ($\pi^e$↑), $C\to D$}

\textbf{a) Demanda.} \emph{Política monetaria}:
\begin{itemize}
\item Expansión: $U<U^n$, $\pi>\pi^e$ ($A\to B$). «Echar leña al fuego»: estimular con la economía en su potencial.
\item \emph{Adaptativas} $\pi^e_t=\pi^e_{t-1}+\gamma(\pi_{t-1}-\pi^e_{t-1})$, $0<\gamma\le1$: la curva sube ($B\to C$); con $\gamma=1$, \concepto{curva aceleracionista}:
\begin{formula}
\centering $\displaystyle \Delta\pi_t\equiv\pi_t-\pi_{t-1}=-a\,(U_t-U^n)$
\end{formula}
\item[] $U<U^n$ exige $\pi$ cada año mayor ($C\to D$): la demanda repetida se vuelve inercia; el menú solo existe a CP. Sostenida, a la larga, monetaria.
\item \emph{Racionales}: islas (\autor{Phelps}, 1970; \autor{Lucas}, 1972) \vertema{3.A.6}: el productor atribuye a causas reales solo la parte $\omega$ (\emph{coeficiente de proyección}); más volatilidad nominal $\Rightarrow$ $\omega$↓, curva casi vertical (\autor{Lucas}, 1973). Con precios rígidos, más vertical cuanto mayor $\pi$ media (\autor{Ball}-\autor{Mankiw}-\autor{Romer}, 1988) \vertema{3.A.7}. \concepto{Inefectividad}: solo la no anticipada mueve $Y$ (\autor{Sargent}-\autor{Wallace}, 1975); \concepto{crítica de \autor{Lucas}}: parámetros cambian con la política.
\item \emph{Sesgo inflacionista} (causa institucional; \autor{Kydland}-\autor{Prescott}, 1977; \autor{Barro}-\autor{Gordon}, 1983): sociedad quiere $U^*<U^n$, banco central discrecional; $\pi$ baja no creíble; en equilibrio $U=U^n$ y \concepto{inflación de discrecionalidad}:
\begin{formula}
\centering $\displaystyle \pi^{\mathrm D}=\pi^*+\frac{\alpha}{\gamma}\,(U^n-U^*)>\pi^*,\quad \alpha=1/a$
\end{formula}
\item[] No compra paro: precio de la falta de compromiso. Remedios: reglas, reputación (juego repetido \vertema{3.A.14}), banco central independiente y conservador (\autor{Rogoff}, 1985), objetivos de inflación (Nueva Zelanda, 1990) \vertema{3.A.36}.
\end{itemize}
\emph{Política fiscal} \vertema{3.A.39}:
\begin{itemize}
\item Restricción del sector público: lo que no cubren impuestos, deuda o dinero; \concepto{señoreaje}: ingreso real por emitir.
\item \emph{Aritmética monetarista desagradable} (\autor{Sargent}-\autor{Wallace}, 1981): fiscal que no estabiliza + banco central que hoy no financia $\Rightarrow$ monetizar más tarde y más; anticipado, $\pi$↑ ya. \concepto{Dominación fiscal}: apretar hoy puede dar más $\pi$ después. Juego de la gallina Tesoro-banco central (\autor{Barthélemy}-\autor{Plantin}, 2019) \vertema{3.A.14}.
\item \emph{Teoría fiscal del nivel de precios} (\autor{Leeper}, 1991; \autor{Sims}, 1994; \autor{Woodford}, 1995): superávits futuros dados $\Rightarrow$ se ajusta $P$ hasta que la deuda real = valor actual de superávits (+ señoreaje):
\begin{formula}
\centering $\displaystyle \frac{(1+i_{t-1})D_{t-1}}{P_tY_t}=\sum_{k\ge0}\Big(\frac{1+\gamma_Y}{1+r}\Big)^{k}(\tau_{t+k}-g_{t+k}+s_{t+k})$
\end{formula}
\item[] Expansión fiscal sin impuestos esperados sube $P$ aunque $M$ no crezca; \autor{Cochrane} (2023), EE.\,UU.\ 2021-22. \emph{Crítica} (\autor{Buiter}, 2002): toma una restricción del Gobierno por condición de equilibrio.
\item \emph{Otras de demanda}: confianza, riqueza, ahorro acumulado, inversión, demanda exterior, depreciación; en NEK suben $r^n_t$: si $i$ no sube bastante, brecha y $\pi$↑. Mismo signo $Y$ y $\pi$ $\Rightarrow$ \concepto{divina coincidencia} \vertema{3.A.7}.
\end{itemize}

\textbf{b) Costes.} Producción ↓ y precios ↑: \concepto{estanflación} (crisis del petróleo de los 70).
\begin{itemize}
\item \emph{NEK} (\autor{Calvo}, 1983: prob.\ $\eta$ de \emph{no} revisar) \vertema{3.A.7}:
\funcionclave{\pi_t=\beta\,\mathbb{E}_t[\pi_{t+1}]+\kappa\,x_t+u_t}
\item[] $u_t$: \concepto{perturbación de costes} (sube coste marginal o margen sin mover la brecha); la misma que $\varepsilon_t$. \emph{Dilema}: con $u_t>0$ no se estabilizan a la vez $\pi$ y $x$ \vertema{3.A.36}: transitoria y expectativas ancladas $\Rightarrow$ «mirar a través»; riesgo de 2.ª ronda $\Rightarrow$ subir tipos.
\end{itemize}
\graficorepaso[\linewidth]{nek-xpi}{a) Demanda: $\pi$ y $x$ juntas, el tipo estabiliza las 2 (divina coincidencia). b) Costes: el dilema}
\begin{itemize}
\item \emph{¿Por qué un precio relativo sube el nivel?} Con costes de cambiar precios, solo cambian los de perturbación grande: pocos suben mucho, muchos esperan $\Rightarrow$ $P$↑ (\autor{Ball}-\autor{Mankiw}, 1995).
\item \emph{3 pasos} (BdE, 2022): \emph{directos} (energía en el IPC); \emph{indirectos} (insumos trasladados); \emph{2.ª ronda}: salarios para recuperar poder adquisitivo + márgenes $\Rightarrow$ \concepto{espiral de precios y salarios}. «La pescadilla que se muerde la cola»: convierte lo transitorio en persistente.
\end{itemize}

\textbf{c) Inercia.} \concepto{Persistencia}: tendencia de $\pi$ a quedarse donde estaba. «Genio y figura, hasta la sepultura».
\begin{itemize}
\item 2 fuentes (\autor{Fuhrer}-\autor{Moore}, 1995): \concepto{heredada} (persisten sus determinantes: paro, brecha, coste marginal); \concepto{intrínseca} ($\pi$ pasada influye por sí: indexación, expectativas que miran atrás, contratos; las $a_i$). 3.ª: retrasos de la política. Si intrínseca, desinflar cuesta aun con credibilidad.
\item \emph{Modelo}: NKPC, $u_t=\rho_u u_{t-1}+$ innovación; hacia delante:
\begin{formula}
\centering $\displaystyle \pi_t=\kappa\sum_{k\ge0}\beta^k\mathbb{E}_t[x_{t+k}]+\frac{u_t}{1-\beta\rho_u}$
\end{formula}
\item[] Sin $\pi$ pasada: toda la persistencia es heredada (la de $\rho_u$). Rigidez en el \emph{nivel}, no en la tasa. \emph{Impl.}: desinflación creíble sin coste (incluso auge: \autor{Ball}, 1994), pero en los datos cuesta; NKPC $\neq$ aceleracionista: con $U<U^n$, $\pi$ alta hoy y se espera que \emph{baje} (\autor{Mankiw}, 2001).
\item \emph{Extensiones}: indexación (\autor{Christiano}-\autor{Eichenbaum}-\autor{Evans}, 2005) o reglas sencillas (\autor{Galí}-\autor{Gertler}, 1999) $\Rightarrow$ \concepto{\autor{Phillips} híbrida} = triángulo con expectativas racionales:
\begin{formula}
\centering $\displaystyle \pi_t=\gamma_b\pi_{t-1}+\gamma_f\mathbb{E}_t[\pi_{t+1}]+\kappa x_t+u_t$
\end{formula}
\item[] Indexación completa: $\gamma_b=\frac{1}{1+\beta}$, $\gamma_f=\frac{\beta}{1+\beta}$ (casi mitad y mitad). Salarios relativos (\autor{Fuhrer}-\autor{Moore}); \emph{información rígida} (\autor{Mankiw}-\autor{Reis}, 2002): inercia sin indexación, desinflación costosa; \emph{expectativas}: ancladas, las perturbaciones pasan; desancladas, se perpetúa.
\item \emph{Evid.}: peso prospectivo > retrospectivo (\autor{Galí}-\autor{Gertler}); persistencia alta sin ancla (70s), casi nula con objetivos creíbles (\autor{Benati}, 2008): crítica de \autor{Lucas}. \emph{Val.}: indexación choca con microdatos (mitad de precios dura >4 meses: \autor{Bils}-\autor{Klenow}, 2004). Pesa más con expectativas mal ancladas y salarios, pensiones, alquileres indexados: desindexación española (2015).
\end{itemize}

\textbf{d) La inflación de 2021-2026.} España: IPC 10,8\,\% (jul.\ 2022; IAPC 10,7); zona del euro 10,6\,\% (oct.\ 2022).
\begin{itemize}
\item \emph{Costes}: cuellos de botella y Ucrania. Zona del euro: energía 3,8 pp (2022), alimentos 2,9, cuellos 1,4 (inicio 2023); mercado laboral solo 0,2 (\autor{Arce} et al., 2024). Beneficios unitarios $\approx$45\,\% del deflactor del consumo (2022T1-2023T1), importación 40\,\% (\autor{Hansen}-\autor{Toscani}-\autor{Zhou}, 2023).
\item \emph{Demanda}: EE.\,UU., estímulo 2021 >> brecha + ahorro (\autor{Summers}, \autor{Blanchard}, 2021); curva plana en los 2010 que se empina con mercado tenso (\autor{Benigno}-\autor{Eggertsson}, 2023); precios cambiados más a menudo.
\item \emph{Inercia}: nació de perturbaciones de precios; la tensión laboral, más tarde, la hizo persistir (\autor{Bernanke}-\autor{Blanchard}, 2025; \autor{Arce} et al.). Expectativas $\approx$2\,\%: bajó sin recesión.
\item \emph{2026}: energía ↑, general ↑, subyacente no: oct., zona del euro 2,5\,\%, España (INE) 3,1\,\%; EE.\,UU.\ 3,7\,\% (sept.). BCE sube en jun.\ y sept.\ (depósito 2,25 y 2,50\,\%) para no desanclar \vertema{3.A.36}.
\end{itemize}
\graficorepaso[0.7\linewidth]{inflacion-general-subyacente}{General y subyacente, zona del euro y España, 1999-2026: en 2022 la subyacente siguió con retraso (2.ª ronda); en 2026, de momento, no}

% ==========================================================================
\bloque[8]{II. Efectos sobre la eficiencia y el bienestar\label{rep:II}}
\begin{ideaclave}
Cuesta aunque se anticipe (encarece el dinero, obliga a cambiar precios, distorsiona impuestos); si no se anticipa, además redistribuye y crea incertidumbre. Óptimo: $-r$ con solo la fricción del dinero, 0 con precios rígidos, algo $>0$ con límite inferior y salarios rígidos: el 2\,\%.
\end{ideaclave}

\sub{II.1. Costes de la anticipada y la no anticipada\label{rep:II.1}}
\begin{itemize}
\item \autor{Fisher}: $i\simeq r+\pi^e$ \vertema{3.A.3}; real efectivo $i-\pi$. \concepto{Anticipada} ($\pi=\pi^e$): nadie gana ni pierde por sorpresa; cuesta protegerse (\emph{eficiencia}). \concepto{No anticipada} ($\pi\neq\pi^e$): redistribuye (\emph{bienestar}). Frontera borrosa: alta = variable = mal anticipada (\autor{Friedman}, 1977).
\item \emph{Suela de zapatos}: tener dinero cuesta $i$; menos saldos, más gestiones \vertema{3.A.35}. \autor{Bailey} (1956): $L(i)\equiv M^D/P$; eficiente $L(0)$; se pierde el área bajo la demanda menos el rectángulo $iL(i)$ (impuesto inflacionario, transferencia al emisor):
\funcionclave{\mathrm{PB}(i)=\int_0^{i}L(x)\,dx-i\,L(i)\approx\tfrac12\lvert L'(0)\rvert\,i^2}
\item[] Crece con el \emph{cuadrado}: del 2 al 4\,\% cuesta 3 veces más que del 0 al 2\,\%. Con \autor{Cagan}, $\lvert L'(0)\rvert=\alpha\zeta$. \emph{Evid.}: EE.\,UU.\ 1900-94: del 10 al 0\,\% $\approx$ algo menos del 1\,\% de renta (\autor{Lucas}, 2000).
\end{itemize}
\graficorepaso[0.62\linewidth]{bailey}{Con $i_0=r+\pi$: rectángulo = impuesto inflacionario (transferencia); triángulo = pérdida de eficiencia, nula con $i=0$}
\begin{itemize}
\item \concepto{Costes de menú}: pequeños, pero bastan para que los precios no se ajusten a la vez (\autor{Mankiw}, 1985); mayor el indirecto: \emph{dispersión de precios relativos} $\Rightarrow$ en NEK $\pi$ pesa mucho en la pérdida; estabilidad = óptimo con rigidez (\autor{Woodford}, 2003) \vertema{3.A.7}.
\item \emph{Unidad de cuenta}: contratos, contabilidad, comparaciones peor. «Las cuentas claras y el chocolate espeso». Volátil: precios señalan peor (\autor{Lucas}, 1973; \autor{Friedman}, 1977).
\item \emph{Fiscales}: \emph{progresividad en frío} (tramos sin actualizar); \emph{rentas del capital}: $(1-\tau)i-\pi=(1-\tau)r-\tau\pi$ ($r=2$, $\tau=30$, $\pi=2$\,\%: 1,4 $\to$ 0,8\,\%); amortizaciones a coste histórico. \autor{Feldstein} (1997): 2\,\% cuesta $\approx$1\,\% PIB \emph{cada año} vs.\ desinflar $\approx$5\,\% una vez. \emph{\autor{Olivera}-\autor{Tanzi}} (1967, 1977): retraso en el cobro; 2 meses al 10\,\% mensual, $-17$\,\% ($1/1{,}1^2\simeq0{,}83$): eslabón de hiperinflaciones.
\item \emph{Capital}: \concepto{\autor{Mundell}-\autor{Tobin}} (1963, 1965): dinero alternativo al capital $\Rightarrow$ $K$↑ \vertema{3.A.43}; \concepto{superneutralidad} (\autor{Sidrauski}, 1967): $K$ por la regla de oro modificada \vertema{3.A.35}; pago por adelantado de la inversión: $K$↓ (\autor{Stockman}, 1981). Datos: más $\pi$, menos inversión y crecimiento (volatilidad).
\item \emph{No anticipada}: \emph{redistribución} de acreedores a deudores (pierden mayores y ricos con bonos; ganan jóvenes con hipoteca fija y sector público: \autor{Doepke}-\autor{Schneider}, 2006); licuación de deuda pública tras la II Guerra Mundial (\autor{Reinhart}-\autor{Sbrancia}, 2015) \vertema{3.A.39}; \emph{regresividad} (más efectivo y energía en la cesta: BCE, 2022); \emph{incertidumbre}: contratos cortos, prima de riesgo de inflación; más $\pi$ puede ir con \emph{más} paro (\autor{Friedman}, 1977); \emph{competitividad}: con cambio fijo o unión, TCR↑ \vertema{3.B.11}.
\item \emph{¿Beneficios?} Lubricante del mercado laboral (\autor{Tobin}, 1972; \autor{Akerlof}-\autor{Dickens}-\autor{Perry}, 1996) \vertema{3.A.28}; señoreaje con fisco pobre (\autor{Phelps}, 1973); margen frente al límite inferior \vertema{3.A.37}. Ninguno justifica $\pi$ alta.
\end{itemize}

\sub{II.2. ¿Existe una tasa de inflación óptima?\label{rep:II.2}}
\begin{itemize}
\item 3 respuestas: deflación, estabilidad o algo de inflación (\autor{Schmitt-Grohé}-\autor{Uribe}, 2010).
\item \concepto{Regla de \autor{Friedman}} (1969): coste privado del dinero $i$ = social $\approx0$. Supuestos: precios flexibles, dinero útil y gratis, impuestos de suma fija. \autor{Sidrauski}: $u_m/u_c=i$; saciar, $u_m=0$:
\funcionclave{i=r+\pi=0\ \Longrightarrow\ \pi^*=-r}
\item[] Con $r$ del 2-4\,\%, $\pi^*$ de $-2$ a $-4$\,\%; nadie la sigue. \emph{Ext.}: impuestos distorsionantes $\Rightarrow$ $\pi^*>-r$ (\autor{Phelps}, 1973); dinero bien intermedio $\Rightarrow$ vuelve $-r$ (\autor{Chari}-\autor{Christiano}-\autor{Kehoe}, 1996); rentas no gravables (harían falta >30\,\%); demanda exterior de moneda (EE.\,UU.\ 2-10\,\%). \emph{Val.}: con precios rígidos, óptimo \emph{cero} (\autor{Woodford}, 2003); con las 2 fricciones, entre $-r$ y 0, cerca de 0.
\item \emph{¿Por qué el 2\,\%?} \emph{Sesgo de medición}: EE.\,UU., 1,1 pp/año (0,6 calidad) (\autor{Boskin}, 1996) \vertema{3.A.9}; \emph{límite inferior}: antes de 2008 riesgo pequeño ($i$ óptimo medio 4,4\,\%, d.t.\ 0,9; $\pi$ óptima $-0{,}4$\,\%), después mayor: se propuso estudiar un 4\,\% (\autor{Blanchard}-\autor{Dell'Ariccia}-\autor{Mauro}, 2010) \vertema{3.A.40}; \emph{salarios rígidos a la baja}: < medio punto. «Curarse en salud». En unión monetaria, evita deflación para abaratarse.
\item \emph{Práctica}: BCE, 2\,\% \emph{simétrico} a medio plazo (jul.\ 2021; jun.\ 2025); Fed, 2\,\% con el deflactor del consumo (inflación media, ago.\ 2020; abandonada ago.\ 2025).
\end{itemize}

\sub{II.3. Inflación y estabilización\label{rep:II.3}}
\begin{itemize}
\item Demanda: estabilizar $\pi$ estabiliza $Y$ (divina coincidencia); costes: elegir según el peso de cada una \vertema{3.A.36}. Objetivo creíble = flexibilidad: con expectativas ancladas se mira a través sin que se vuelva inercia (2022-24; a prueba en 2026).
\end{itemize}

% ==========================================================================
\bloque[5,5]{III. Contextos extremos: hiperinflación y deflación\label{rep:III}}
\begin{ideaclave}
La hiperinflación es fiscal en su origen y monetaria en su mecanismo, y se acaba de golpe con una reforma fiscal creíble; la deflación es buena si viene de la oferta y peligrosa si viene de una demanda hundida con deudas altas.
\end{ideaclave}

\sub{III.1. Hiperinflación\label{rep:III.1}}
\begin{itemize}
\item \emph{Idea}: $>50$\,\% al mes $\approx$13.000\,\% al año ($1{,}5^{12}-1\simeq128{,}7$). 56 hasta 2012 (guerras, revoluciones, caída de regímenes); Venezuela la 57.ª (\autor{Hanke}-\autor{Krus}, 2013). 2 cosas: Estado sin impuestos ni deuda que recurre al banco central + pérdida de confianza en el dinero.
\item \emph{Máximos mensuales}: Hungría jul.\ 1946 $4{,}19\cdot10^{16}$\,\% (×2 en 15 h); Zimbabue nov.\ 2008 (24,7 h); Yugoslavia ene.\ 1994; Alemania oct.\ 1923 29.500\,\% (3,7 días); Grecia 1944; China 1949. Venezuela desde nov.\ 2016: 130.060\,\% anual (2018).
\item \concepto{Señoreaje} e \concepto{impuesto inflacionario} \vertema{3.A.39}:
\begin{formula}
\centering $\displaystyle S_t\equiv\frac{M_t-M_{t-1}}{P_t}=\Big(\frac{M_t}{P_t}-\frac{M_{t-1}}{P_{t-1}}\Big)+\frac{\pi_t}{1+\pi_t}\frac{M_{t-1}}{P_{t-1}}$
\end{formula}
\item[] Base: saldos reales; tipo: $\pi/(1+\pi)$; estacionario continuo, $S=\pi M/P$.
\item \emph{\autor{Cagan} (1956)}: renta y $r$ apenas cambian; demanda $\frac{M^D_t}{P_t}=\zeta e^{-\alpha\pi^e_t}$ ($\alpha$ semielasticidad); adaptativas $\dot\pi^e=\gamma(\pi-\pi^e)$. Estacionario $\pi=\pi^e=\gamma_M$:
\begin{formula}
\centering $\displaystyle S(\pi)=\pi\zeta e^{-\alpha\pi};\quad S'=0\iff\pi=\tfrac1\alpha;\quad S^{\max}=\tfrac{\zeta}{\alpha e}$
\end{formula}
\item[] \concepto{\autor{Laffer} del señoreaje}: déficit $<S^{\max}$, 2 inflaciones ($\pi_B$, $\pi_A$); mayor, ninguna. \emph{Estabilidad}:
\funcionclave{\dot\pi^e=\frac{\gamma}{1-\alpha\gamma}(\gamma_M-\pi^e):\ \text{estable si }\alpha\gamma<1}
\item[] $\alpha\gamma>1$: la huida del dinero se alimenta sola («quien juega con fuego…»). 2 caminos: déficit $>S^{\max}$ o $\alpha\gamma>1$.
\end{itemize}
\pareja{laffer-senoreaje}{Máximo en $\pi=1/\alpha$; $S$ con $\pi_B$ o $\pi_A$; $S'>S^{\max}$, ninguna}{cagan-estabilidad}{a) $\alpha\gamma<1$, estable; b) $\alpha\gamma>1$, inestable}
\begin{itemize}
\item \emph{Ext.}: racionales: $P$ depende de la senda \emph{esperada} de $M$; anuncio creíble baja $\pi$ ya (\autor{Sargent}-\autor{Wallace}, 1973); círculo fiscal: \autor{Olivera}-\autor{Tanzi} + huida a la divisa.
\item \emph{Evid.}: \autor{Cagan}, 7 hiperinflaciones (Austria, Alemania, Hungría ×2, Polonia, Rusia, Grecia): crecimiento continuo de $M$, demanda que ajusta bien. \concepto{Paradoja de \autor{Cagan}}: $\pi>1/\alpha$; la explica la lentitud de las expectativas.
\item \emph{Fin} (\autor{Sargent}, 1982): Austria, Hungría, Polonia, Alemania (años 20): de golpe al cambiar el régimen (presupuesto equilibrado + banco central que no financia); costes pequeños; $M$ siguió creciendo (remonetización). Fiscal en origen, se apaga con credibilidad.
\end{itemize}

\sub{III.2. Deflación\label{rep:III.2}}
\begin{itemize}
\item \emph{Buena} (oferta: productividad, s.~\textsc{xix}) y \emph{mala} (demanda: Gran Depresión) (\autor{Bordo}-\autor{Filardo}, 2005). Caída puntual por energía no es deflación (2020: tasas negativas ago.-dic., pero año +0,3\,\%).
\item \emph{Evid.}: Japón 1998-2012: IPC $-3{,}9$\,\%, deflactor $-15{,}6$\,\%, tipos $\approx0$; zona del euro: mín.\ $-0{,}6$\,\% (ene.\ 2015), media 0,2\,\% en 2015 y 2016; España $-0{,}2$ / $-0{,}6$ / $-0{,}3$\,\% (2014-16); BCE: depósito negativo y compras de deuda (2015) \vertema{3.A.40}.
\item \emph{Peligros}: \concepto{deflación por deudas} (\autor{Fisher}, 1933): «cuanto más pagan los deudores, más deben» \vertema{3.A.42}; salarios rígidos (real↑, paro↑); aplazar compras; \emph{límite inferior}: $r=i-\pi^e$ \emph{sube} justo cuando hace falta que baje \vertema{3.A.37}.
\item \emph{Trampa} (\autor{Krugman}, 1998, Japón): $r^n<0$ e $i$ en el límite: solo bajan $r$ las expectativas de $\pi$; comprometerse creíblemente a ser irresponsable \vertema{3.A.40}. \autor{Eggertsson}-\autor{Woodford} (2003): tipos bajos más tiempo, como objetivo de nivel de precios. \emph{Reflacionar} (\autor{Fisher}): orientación, compras, fiscal no financiada con impuestos (\autor{Sims}, 2016), helicóptero (\autor{Bernanke}, 2002); BoJ 2013, 2\,\% y compras masivas.
\item \emph{Desinflar} cuesta (inercia): \concepto{ratio de sacrificio}:
\begin{formula}
\centering $\displaystyle \text{RS}=\frac{\sum_t(y^n_t-y_t)}{\pi_{\text{inicial}}-\pi_{\text{final}}}$
\end{formula}
\item[] \autor{Ball} (1994), 65 desinflaciones OCDE: menor si rápida y salarios flexibles; teoría: y si creíble; histéresis la encarece \vertema{3.A.28}. \autor{Volcker}: 14,8\,\% (mar.\ 1980) $\to$ 3,8\,\% (dic.\ 1982), paro 10,8\,\% (nov.\ 1982); 2022-24, sin recesión. Mejor errar algo por arriba.
\end{itemize}

% ==========================================================================
\bloque[1]{Conclusión}
\begin{itemize}
\item \textbf{Recapitulación}: «En conclusión,» a LP monetaria; a CP demanda, costes e inercia (\autor{Gordon}); fiscal y credibilidad deciden si el dinero la alimenta. Costes pequeños si baja y anticipada, grandes si alta y volátil; \autor{Friedman}, óptimo sin rigideces; rigideces, límite inferior y medición $\Rightarrow$ 2\,\%. Hiperinflación: déficit que solo paga la imprenta, acaba con reforma fiscal creíble; deflación de demanda: deudas y monetaria atada.
\item \textbf{Relevancia}: 2021-23 y 2026: los costes se vuelven persistentes si se desanclan expectativas; costes desiguales.
\item \textbf{Extensiones}: \vertema{3.A.35}, \vertema{3.A.36}, \vertema{3.A.37}, \vertema{3.A.39}, \vertema{3.A.40}, \vertema{3.A.7}, \vertema{3.A.28}, \vertema{3.A.9}.
\item \textbf{Opinión}: 2\,\% simétrico, compromiso razonable; lo sostienen credibilidad y fiscal sostenible, no solo el tipo.
\item \textbf{Idea final}: quien juega con fuego se quema, pero una casa sin fuego se hiela: mantener la lumbre baja y encendida.
\end{itemize}

% ==========================================================================
\bloque{Autores y fechas}
\hito{1926}{\autor{Fisher}: inflación y empleo («I Discovered the Phillips Curve», 1973)}
\hito{1933}{\autor{Fisher}: deflación por deudas}
\hito{1956}{\autor{Cagan}: hiperinflación; \autor{Bailey}: coste de bienestar}
\hito{1958}{\autor{Phillips}: salarios y paro (R.\,U.)}
\hito{1960}{\autor{Samuelson}-\autor{Solow}: el menú; \autor{Lipsey}}
\hito{1963}{\autor{Friedman}: «siempre y en todas partes»; \autor{Mundell}}
\hito{1967}{\autor{Phelps}; \autor{Sidrauski}; \autor{Olivera}}
\hito{1968}{\autor{Friedman}: curva aumentada}
\hito{1969}{\autor{Friedman}: la cantidad óptima de dinero}
\hito{1972}{\autor{Lucas}: islas y expectativas racionales}
\hito{1973}{\autor{Phelps}: inflación como impuesto}
\hito{1975}{\autor{Sargent}-\autor{Wallace}: inefectividad}
\hito{1977}{\autor{Kydland}-\autor{Prescott}; \autor{Tanzi}; \autor{Friedman} (Nobel)}
\hito{1981}{\autor{Sargent}-\autor{Wallace}: aritmética desagradable}
\hito{1982}{\autor{Gordon}: triángulo; \autor{Sargent}: fin de 4 hiperinflaciones}
\hito{1983}{\autor{Barro}-\autor{Gordon}; \autor{Calvo}}
\hito{1985}{\autor{Rogoff}: banquero conservador; \autor{Mankiw}: costes de menú}
\hito{1990}{Nueva Zelanda: objetivo de inflación}
\hito{1991}{\autor{Leeper}: monetaria activa/pasiva (teoría fiscal)}
\hito{1994}{\autor{Ball}: ratio de sacrificio; \autor{Sims}: teoría fiscal}
\hito{1995}{\autor{Fuhrer}-\autor{Moore}; \autor{Ball}-\autor{Mankiw}; \autor{McCandless}-\autor{Weber}}
\hito{1996}{\autor{Boskin}; \autor{Akerlof}-\autor{Dickens}-\autor{Perry}}
\hito{1997}{\autor{Feldstein}: costes fiscales}
\hito{1998}{\autor{Krugman}: trampa en Japón}
\hito{1999}{\autor{Galí}-\autor{Gertler}: curva híbrida}
\hito{2000}{\autor{Lucas}: coste de bienestar}
\hito{2002}{\autor{Mankiw}-\autor{Reis}: información rígida}
\hito{2003}{\autor{Woodford}; \autor{Eggertsson}-\autor{Woodford}}
\hito{2005}{\autor{Christiano}-\autor{Eichenbaum}-\autor{Evans}}
\hito{2010}{\autor{Schmitt-Grohé}-\autor{Uribe}: tasa óptima}
\hito{2021}{BCE: 2\,\% simétrico (confirmado 2025)}
\hito{2025}{\autor{Bernanke}-\autor{Blanchard}: inflación 2021-23}

% ==========================================================================
\bloque{Datos clave}
\begin{itemize}
\item Inflación oct.\ 2026: zona del euro \dato{4,0\,\%}{Eurostat}; energía \dato{+20,8\,\%}{Eurostat}; España \dato{4,9\,\%}{INE}; subyacente \dato{2,5\,\%}{Eurostat} y \dato{3,1\,\%}{INE}; EE.\,UU.\ \dato{3,7\,\%}{BLS, sept.}
\item Expectativas a LP: \dato{2,0\,\%}{BCE, SPF 4T 2026}. Depósito BCE \dato{2,50\,\%}{BCE, sept.\ 2026}.
\item Máximos 2022: España \dato{10,8\,\%}{INE, jul.}; zona del euro \dato{10,6\,\%}{Eurostat, oct.}
\item Media zona del euro 1999-2020: \dato{1,6\,\%}{Eurostat}.
\item Alemania: \dato{29.500\,\% al mes}{Hanke-Krus, oct.\ 1923}; Hungría \dato{$4{,}19\cdot10^{16}$\,\%}{jul.\ 1946}; Venezuela \dato{130.060\,\%}{FMI, 2018}.
\item Japón 1998-2012: IPC \dato{$-3{,}9$\,\%}{BM}; deflactor \dato{$-15{,}6$\,\%}{BM}.
\item \autor{Volcker}: \dato{14,8 $\to$ 3,8\,\%}{BLS, 1980-82}; paro \dato{10,8\,\%}{BLS, nov.\ 1982}.
\end{itemize}

% ==========================================================================
\bloque{Preguntas del tribunal}
\begin{enumerate}
\item \textbf{¿Es la inflación siempre un fenómeno monetario?} A LP sí (correlación $\approx$0,95); a CP la mueven demanda, costes e inercia, pero solo se \emph{sostiene} si el dinero (o el déficit) la alimenta.
\item \textbf{¿Qué diferencia la curva aceleracionista de la NKPC?} Aceleracionista: $U<U^n$ $\Rightarrow$ $\pi$ cada vez mayor; NKPC: $\pi$ alta hoy y se espera que baje; toda su persistencia es heredada.
\item \textbf{¿Por qué la pérdida de \autor{Bailey} crece con el cuadrado de $i$?} Es un exceso de gravamen: $\mathrm{PB}\approx\frac12\lvert L'(0)\rvert i^2$; del 2 al 4\,\% cuesta el triple que del 0 al 2\,\%.
\item \textbf{¿Si la regla de \autor{Friedman} es óptima, por qué el 2\,\%?} Rigideces (óptimo 0), límite inferior, salarios rígidos a la baja y sesgo de medición; en una unión monetaria, evitar deflación al ajustar.
\item \textbf{¿Por qué se dispara una hiperinflación si más dinero no da más señoreaje?} Déficit $>S^{\max}$ o $\alpha\gamma>1$ (huida que se alimenta); paradoja de \autor{Cagan}: acelerar paga mientras las expectativas tardan.
\item \textbf{¿Cómo se acaba?} De golpe, con cambio de régimen fiscal y monetario creíble (\autor{Sargent}, 1982); el dinero sigue creciendo por la remonetización.
\item \textbf{¿Es mala toda deflación?} No: la de oferta acompaña al crecimiento; la de demanda con deudas altas y tipos en el límite sube la carga real y el tipo real.
\end{enumerate}

\begin{trampa}
\begin{itemize}
\item Hiperinflación (\autor{Cagan}): $>50$\,\% \emph{al mes}, no al año ($\approx$13.000\,\% anual).
\item Deflación: caída \emph{sostenida}; un mes negativo no basta. Desinflación: $\pi$↓ pero $>0$.
\item Regla de \autor{Friedman} (1969): $\pi^*=-r$, $i=0$; no es la regla del $k$\,\% (1960) ni el helicóptero.
\item Señoreaje $\neq$ impuesto inflacionario: coinciden solo si los saldos reales no cambian.
\item Máximo del señoreaje en $\pi=1/\alpha$, $S^{\max}=\zeta/(\alpha e)$; estable si $\alpha\gamma<1$.
\item \autor{Calvo}: $\eta$ = prob.\ de \emph{no} revisar; más $\eta$, NKPC más plana.
\item Triángulo: pesos $\sum a_i=1$ (aceleracionista); la suma empieza en $i=1$.
\item \autor{Mundell}-\autor{Tobin}: $K$↑; \autor{Sidrauski}: superneutral; \autor{Stockman}: $K$↓.
\item Rentas del capital: rendimiento real neto $(1-\tau)r-\tau\pi$: la inflación lo reduce en $\tau\pi$.
\item Redistribución: de acreedores a deudores (pierden mayores con bonos).
\item Dominación fiscal: apretar hoy puede dar \emph{más} inflación después (\autor{Sargent}-\autor{Wallace}).
\item Divina coincidencia: con perturbaciones de demanda, no de costes.
\end{itemize}
\end{trampa}

\end{columnas}
\end{document}
